\documentclass[aspectratio=169]{beamer}

% --- Tiếng Việt ---
\usepackage[utf8]{inputenc}
\usepackage[vietnamese]{babel}

% --- Toán ---
\usepackage{amsmath, amssymb}

% --- Vẽ hình ---
\usepackage{tikz}
\usetikzlibrary{arrows.meta, positioning}

% --- Giao diện ---
\usetheme{Madrid}
\usecolortheme{seahorse}

% --- Xóa footer ---
\setbeamertemplate{footline}{}
\setbeamertemplate{navigation symbols}{}

% --- Màu ---
\definecolor{main}{RGB}{30,60,150}

% =====================
\title{}
\author{}
\date{}
% =====================

\begin{document}

% =====================================
% SLIDE 1: TRANG BÌA (RẤT ĐẸP)
% =====================================
\begin{frame}
\begin{tikzpicture}[remember picture, overlay]
\fill[blue!15] (current page.south west) rectangle (current page.north east);
\fill[blue!70] (current page.south west) rectangle ([yshift=3cm]current page.south east);
\end{tikzpicture}

\vspace{2.5cm}

\begin{center}
{\Huge \textbf{LÝ THUYẾT ĐỒ THỊ}}\\[0.8cm]

{\LARGE \textbf{BÀI 1: CÁC KHÁI NIỆM CƠ BẢN}}\\[0.5cm]

\rule{9cm}{1.2pt}
\end{center}

\vfill
\begin{flushright}
\textit{Môn: Tin học}\\
\textit{Chủ đề: Graph}
\end{flushright}

\end{frame}

% =====================================
% SLIDE 2: GIỚI THIỆU
% =====================================
\begin{frame}{Giới thiệu}
\begin{itemize}
    \item Đồ thị là một trong những cấu trúc dữ liệu quan trọng nhất trong Tin học.
    \item Được sử dụng để mô tả mối quan hệ giữa các đối tượng.
    \item Ứng dụng rộng rãi trong:
    \begin{itemize}
        \item Mạng xã hội
        \item Bản đồ và định vị GPS
        \item Trí tuệ nhân tạo
    \end{itemize}
\end{itemize}

\begin{block}{Mục tiêu bài học}
Hiểu rõ định nghĩa và các khái niệm cơ bản của đồ thị.
\end{block}
\end{frame}

% =====================================
% I. ĐỊNH NGHĨA ĐỒ THỊ
% =====================================
\begin{frame}{I. Định nghĩa đồ thị}
\begin{block}{Khái niệm}
Đồ thị là một cấu trúc gồm hai thành phần:
\[
G = (V, E)
\]
\end{block}

\begin{itemize}
    \item $V$: tập hợp các đỉnh (vertices)
    \item $E$: tập hợp các cạnh (edges)
\end{itemize}

\pause

\begin{block}{Ý nghĩa}
Đồ thị dùng để biểu diễn các mối quan hệ giữa các đối tượng trong thực tế.
\end{block}
\end{frame}

% =====================================
% SLIDE: HÌNH MINH HỌA
% =====================================
\begin{frame}{Ví dụ minh họa đồ thị}
\begin{center}
\begin{tikzpicture}[
    every node/.style={circle, draw, fill=blue!20, minimum size=1cm},
    node distance=2.5cm
]

\node (A) {A};
\node (B) [right of=A] {B};
\node (C) [below of=A] {C};
\node (D) [below of=B] {D};

\draw[thick] (A)--(B);
\draw[thick] (A)--(C);
\draw[thick] (B)--(D);
\draw[thick] (C)--(D);

\end{tikzpicture}
\end{center}

\begin{block}{Mô tả}
\begin{itemize}
    \item $V = \{A, B, C, D\}$
    \item $E = \{AB, AC, BD, CD\}$
\end{itemize}
\end{block}
\end{frame}

% =====================================
% II. CÁC KHÁI NIỆM
% =====================================
\begin{frame}{II. Các khái niệm cơ bản}
\begin{itemize}
    \item \textbf{Đỉnh (Vertex)}: điểm biểu diễn đối tượng
    \item \textbf{Cạnh (Edge)}: liên kết giữa hai đỉnh
    \item \textbf{Đường đi}: dãy các cạnh liên tiếp
\end{itemize}

\pause

\begin{block}{Ví dụ}
A → B → D là một đường đi trong đồ thị.
\end{block}
\end{frame}

% =====================================
% SLIDE: BẬC ĐỈNH
% =====================================
\begin{frame}{Bậc của đỉnh}
\begin{block}{Định nghĩa}
Bậc của một đỉnh là số cạnh nối với đỉnh đó.
\end{block}

\begin{itemize}
    \item deg(A) = 2
    \item deg(B) = 2
    \item deg(C) = 2
\end{itemize}

\pause

\begin{block}{Ý nghĩa}
Cho biết mức độ liên kết của một đỉnh trong đồ thị.
\end{block}
\end{frame}

% =====================================
% SLIDE: PHÂN LOẠI
% =====================================
\begin{frame}{Phân loại đồ thị}
\begin{itemize}
    \item \textbf{Đồ thị vô hướng}: cạnh không có chiều
    \item \textbf{Đồ thị có hướng}: cạnh có hướng
\end{itemize}

\pause

\begin{block}{Ví dụ}
\begin{itemize}
    \item Facebook → vô hướng
    \item Google Maps → có hướng
\end{itemize}
\end{block}
\end{frame}

% =====================================
% SLIDE: KẾT LUẬN
% =====================================
\begin{frame}{Kết luận}
\begin{itemize}
    \item Đồ thị là công cụ quan trọng trong Tin học
    \item Giúp mô hình hóa các bài toán thực tế
\end{itemize}

\begin{block}{Ghi nhớ}
Nắm chắc khái niệm là nền tảng để học BFS, DFS.
\end{block}
\end{frame}

% =====================================
% SLIDE: END
% =====================================
\begin{frame}
\begin{center}
\Huge \textbf{Cảm ơn đã lắng nghe!}
\end{center}
\end{frame}

\end{document}